\documentclass[11pt]{article}
\usepackage[T1]{fontenc}
\usepackage[utf8]{inputenc}
\usepackage{lmodern,microtype}
\usepackage[a4paper,margin=28mm]{geometry}
\usepackage{amsmath,amssymb,amsthm,mathtools}
\usepackage{enumitem}
\usepackage[hidelinks]{hyperref}
\usepackage[nameinlink,noabbrev]{cleveref}
\hypersetup{pdftitle={A nonfree compact projective condensed abelian group},pdfsubject={A mathematical note on Boolean retractions and condensed abelian groups}}
\newtheorem{theorem}{Theorem}[section]
\newtheorem{proposition}[theorem]{Proposition}
\newtheorem{lemma}[theorem]{Lemma}
\newtheorem{corollary}[theorem]{Corollary}
\theoremstyle{definition}
\newtheorem{definition}[theorem]{Definition}
\newtheorem{convention}[theorem]{Convention}
\theoremstyle{remark}
\newtheorem{remark}[theorem]{Remark}
\newcommand{\Z}{\mathbb Z}
\newcommand{\N}{\mathbb N}
\newcommand{\CAb}{\operatorname{Cond}(\mathrm{Ab})}
\newcommand{\Clop}{\operatorname{Clop}}
\newcommand{\Stone}{\operatorname{Stone}}
\newcommand{\RO}{\operatorname{RO}}
\newcommand{\Res}{\operatorname{Res}}
\newcommand{\Eq}{\operatorname{Eq}}
\newcommand{\BA}{\operatorname{BA}_{\operatorname{tr}(3)}}
\newcommand{\restr}{\mathbin{\upharpoonright}}
\newcommand{\cc}{\mathfrak c}
\newcommand{\tr}{\operatorname{tr}}
\newcommand{\id}{\operatorname{id}}
\setlist[enumerate]{label=\textup{(\roman*)},leftmargin=*,itemsep=3pt}
\setlength{\emergencystretch}{2em}
\title{A nonfree compact projective condensed abelian group}
\author{}
\date{19 October 2026}
\begin{document}
\maketitle
% Begin introduction.tex
\begin{abstract}
We give a negative answer to Question~3.6 of Clausen and Scholze's
\emph{Condensed Mathematics and Complex Geometry}: a compact projective
condensed abelian group need not be free on an extremally disconnected
compact Hausdorff space. We construct such a group as the kernel of a
split map $\Z[S]\to\Z[T]$, where $T\subseteq S$ is a nowhere dense
retract and both spaces are compact Hausdorff and extremally disconnected.
The obstruction to freeness is topological:
no nonempty clopen of $S$ disjoint from $T$ maps onto a nonempty clopen
of $T$. We obtain the required spaces by Stone duality from complete
ccc Boolean algebras, using Shelah's full-super tree families and a
recursion of Boolean free products with a fixed complete target.
\end{abstract}

\section{Introduction}
\label{sec:introduction}

For a compact Hausdorff extremally disconnected space $X$, let $\Z[X]$
denote the free condensed abelian group on its associated condensed set.
These groups are compact projective, and every compact projective
condensed abelian group is a direct summand of one of them
\cite[Proposition~3.5]{ClausenScholzeComplex}.
Clausen and Scholze ask whether every such summand is again of this
form \cite[Question~3.6]{ClausenScholzeComplex}. Here compactness of an
object means that its Hom functor preserves filtered colimits.

\begin{theorem}[A negative answer to Question~3.6]
\label{thm:main}
There is a nonzero compact projective condensed abelian group $P$ such
that $P\not\cong\Z[X]$ for every compact Hausdorff extremally disconnected
space $X$. More precisely, there are such spaces $S,T$ and continuous maps
\[
 r:S\longrightarrow T,\qquad i:T\longrightarrow S,
 \qquad r i=\id_T,
\]
for which
\[
 P=\ker\bigl(\Z[r]:\Z[S]\longrightarrow\Z[T]\bigr)
\]
has this property. In particular,
$\Z[S]\cong\Z[T]\oplus P$.
\end{theorem}

The proof starts with a topological obstruction to freeness. We arrange
that $i(T)$ is nowhere dense in $S$ and that no nonempty clopen of $S$
disjoint from $i(T)$ admits a continuous surjection onto a nonempty
clopen of $T$. Every morphism between free condensed groups on compact
Hausdorff extremally disconnected spaces has a description by finitely
many continuous branches on clopen domains.
Applied to this obstruction, the description rules out a nonzero
scalar multiple of $\id_{\Z[T]}$ factoring through any finite power
of $P$. If $P$ were $\Z[X]$, the branches of its inclusion into $\Z[S]$
would give precisely such a factorization. This proves the criterion
in \cref{prop:transfer}.

We construct the spaces through Boolean algebras. A Boolean retraction
$p:B\to D\subseteq B$ gives the maps $r,i$ by Stone duality. An order
dense kernel makes $i(T)$ nowhere dense; the required obstruction says
that no positive interval of $D$ embeds into an interval of $B$ whose
unit is in $\ker p$. Shelah's full-super tree families supply algebras
$E_\gamma$ that are excluded from completions represented over the
other trees (\cref{prop:exclusion}). Local free-product insertions
preserve these representations on complementary intervals
(\cref{lem:complementary-omission}).

The final construction fixes a complete Boolean algebra $Q$ in which
every positive interval contains every $E_\gamma$. We recursively
construct a complete ccc algebra $B$ and a homomorphism $q:B\to Q$.
Each positive element sent to zero is treated once: an insertion on
its complement excludes a suitable $E_\gamma$ from its interval.
The inserted factors give a lift $H\to B$ of a dense subalgebra
$H\subseteq Q$. Extending this lift gives $s:Q\to B$; since $qs$
fixes $H$, density forces $qs=\id_Q$. Taking $D=s(Q)$ and $p=sq$ gives
the Boolean algebras needed for \cref{prop:transfer}, and completes
the proof of \cref{thm:main} in \cref{sec:construction}.

\paragraph{Boolean conventions.}
For a Boolean algebra $B$ and $b\in B$, write
$B\restr b=\{x\in B:x\leq b\}$, with relative unit $b$ and relative
complement $x\mapsto b\wedge\neg x$. Homomorphisms between intervals
preserve their relative units. All Boolean homomorphisms preserve
finite Boolean operations; an embedding is called \emph{complete}
if it also preserves arbitrary joins. Completeness of a subalgebra
$D\subseteq B$ refers to joins computed in $D$ and does not imply
that its inclusion in $B$ is complete. An algebra is \emph{ccc} if
every disjoint family of nonzero elements is countable. An ideal is
\emph{order dense} if it contains a nonzero element below each
nonzero element of the algebra.

The external set-theoretic inputs are Shelah's full-super existence
theorem and the stated chain-condition and regularity results for
tree Boolean algebras. We specify their versions and hypotheses at
their uses. The parameter-transfer argument, the representations of
completed intervals, the Boolean recursion, and the condensed
implication are proved below.

% End introduction.tex
\tableofcontents
% Begin transfer.tex
\section{A topological obstruction to freeness}
\label{sec:transfer}

\subsection{Free condensed groups on extremally disconnected spaces}

We use the site of profinite sets with finite jointly surjective
covers, or equivalently its extremally disconnected presentation.
The condensed set associated to a compact Hausdorff space $Y$ assigns
$\operatorname{Cont}(U,Y)$ to a test space $U$. Its free condensed
abelian group $\Z[Y]$ is the sheafification of
\[
 U\longmapsto\Z[\operatorname{Cont}(U,Y)].
\]
These conventions and the extremally disconnected presentation are
described in \cite[Theorem~2.2 and Proposition~2.7]{ScholzeCondensed}.

For size conventions, choose an uncountable strong-limit cardinal
$\kappa$ larger than the underlying sets of the spaces in a given
construction, and use $\kappa$-small test spaces. Such bounds exist in
ZFC. The enlargement functors between these categories are fully
faithful and compatible with the free objects and split diagrams used
here \cite[Proposition~2.9, Remark~2.10, Definition~2.11, and
Remark~2.13]{ScholzeCondensed}. A proposed parameter $X$ may therefore
be included by increasing the bound. This makes the final assertion
independent of its initial choice.

\begin{lemma}
\label{lem:finite-branches}
Let $X,Y$ be compact Hausdorff extremally disconnected spaces.
\begin{enumerate}
\item The object $\Z[X]$ is compact projective, and $\Z[Y](*)$ is the
free abelian group on the underlying set of $Y$.
\item Every morphism $\varphi:\Z[X]\to\Z[Y]$ has a finite clopen
partition $X=\coprod_j X_j$ such that, on each $X_j$, its generator
map is of the form
\[
 x\longmapsto\sum_{a=1}^{m_j}n_{j,a}[f_{j,a}(x)],
 \qquad n_{j,a}\in\Z,
\]
where each $f_{j,a}:X_j\to Y$ is continuous. The empty sum is allowed.
\end{enumerate}
\end{lemma}

\begin{proof}
By the free-object adjunction,
$\operatorname{Hom}(\Z[X],M)=M(X)$. In the extremally disconnected
presentation, evaluation at $X$ is exact and preserves filtered
colimits, proving compact projectivity
\cite[proof of Theorem~2.2 and Proposition~2.7, p.~12]{ScholzeCondensed}.

A section of the sheafification over $X$ is locally a section of the
displayed presheaf. More explicitly, the two plus constructions used
in sheafification give finite covers whose composites are still finite
covers; on a refinement the section is represented by finite integral
linear combinations of continuous maps to $Y$. The finite disjoint
union of the covering spaces maps continuously and surjectively to
$X$. A compact Hausdorff extremally disconnected space is projective
with respect to compact Hausdorff surjections, so this map has a
continuous section. Pulling the local representatives back along it
gives a finite clopen partition of $X$ with a finite word on each part.
This proves (ii).

For a point, the identity cover refines every cover: choose a point in
one nonempty covering member. Thus both local representatives and
local equalities in the sheafification pull back to the point. This
shows that sheafification neither adds elements nor introduces
relations in the free abelian group on $Y$ at $*$, proving the
remaining assertion in (i).
\end{proof}

\subsection{The Boolean criterion}

\begin{proposition}
\label{prop:transfer}
Let $D\subseteq B$ be nonzero complete Boolean algebras, and let
$p:B\to D$ be a Boolean retraction. Suppose that $K=\ker p$ is order
dense and that
\begin{equation}
\label{eq:interval-obstruction}
 0<b\in K,\quad 0<d\in D
 \quad\Longrightarrow\quad
 D\restr d\not\hookrightarrow B\restr b.
\end{equation}
Then there is a nonzero compact projective condensed abelian group
which is not isomorphic to $\Z[X]$ for any compact Hausdorff
extremally disconnected space $X$.
\end{proposition}

Let $S=\Stone(B)$ and $T=\Stone(D)$. Stone duality sends the inclusion
and the retraction to maps
\[
 r:S\longrightarrow T,\qquad i:T\longrightarrow S,
 \qquad r i=\id_T.
\]
The map $r$ is surjective and $i$ is a closed embedding. Completeness
of the Boolean algebras makes both spaces extremally disconnected.
We identify $T$ with its image in $S$. A clopen of $S$ is disjoint
from $T$ precisely when its Boolean element lies in $K$. Thus order
density of $K$ says that $T$ has empty interior, and hence is nowhere
dense. Condition \eqref{eq:interval-obstruction} says that there is
no continuous surjection
\begin{equation}
\label{eq:topological-obstruction}
 C\longrightarrow V
\end{equation}
when $C$ is a nonempty clopen of $S$ disjoint from $T$ and $V$ is a
nonempty clopen of $T$.

Define
\begin{equation}
\label{eq:projective-kernel}
 \begin{aligned}
 P&=\ker\bigl(\Z[r]:\Z[S]\to\Z[T]\bigr),\\
 e&=\id_{\Z[S]}-\Z[i]\Z[r].
 \end{aligned}
\end{equation}
Then $e^2=e$ and $\operatorname{im}e=P$. In particular $P$ is compact
projective. It is nonzero: for $s\in S\setminus T$ the element
$e[s]=[s]-[i(r(s))]$ is nonzero in the free group $\Z[S](*)$.

We prove the proposition by establishing two assertions. The interval
condition rules out a nonzero scalar multiple of $\id_{\Z[T]}$
factoring through a finite power of $P$. On the other hand, if
$P\cong\Z[X]$, the finite branches of its inclusion into $\Z[S]$
force just such a factorization, with scalar $1$.

\begin{lemma}
\label{lem:no-scalar-factorization}
For every positive integer $k$ and every nonzero integer $n$, there
are no morphisms
\[
 a:\Z[T]\longrightarrow P^k,\qquad
 b:P^k\longrightarrow\Z[T],\qquad ba=n\id_{\Z[T]}.
\]
\end{lemma}

\begin{proof}
Suppose that $a,b$ exist. Regard $a$ as a map into $\Z[S]^k$ and extend
$b$ to
\[
 q=b\circ\operatorname{diag}(e,\ldots,e):
       \Z[S]^k\longrightarrow\Z[T].
\]
Then $qa=n\id_{\Z[T]}$. For each summand of $\Z[S]^k$, the value of
its $q$-column at every point of $T\subseteq S$ is zero, since
$e[i(t)]=0$.

Apply \cref{lem:finite-branches} to the components of $a$ and $q$.
Denote their incoming branches by $f_\alpha$ and their outgoing
branches by $g_\beta$, retaining the label of the intervening copy
of $S$. At any $t\in T$, the coefficient of $[t]$ in $qa[t]$ equals
$n\ne0$. Group the terms in this composition by their intermediate
point and copy of $S$. Every group passing through a point of $T$
has total contribution zero, because its entire $q$-column is zero.
Consequently there is a compatible pair of branches for which
\begin{equation}
\label{eq:branch-return}
 f_\alpha(t)\notin T,
 \qquad g_\beta(f_\alpha(t))=t.
\end{equation}

There are finitely many branch pairs. For each compatible pair, take
the subset of $T$ on which both branches are defined and
\eqref{eq:branch-return} holds. Its domain conditions are clopen,
the avoidance of $T$ is open, and the equality is closed because
$T$ is Hausdorff. Thus it is locally closed. These sets cover $T$.
A locally closed set with empty interior is nowhere dense; finitely
many such sets cannot cover the nonempty compact Hausdorff space $T$.
One branch pair therefore satisfies \eqref{eq:branch-return} on a
nonempty clopen $V\subseteq T$.

The compact image $f_\alpha(V)$ is disjoint from the closed subset
$T\subseteq S$ and lies in the clopen domain of $g_\beta$. Choose a
clopen $C_0$ of the relevant copy of $S$ containing this image,
disjoint from $T$, and contained in that domain. Set
\[
 C=C_0\cap g_\beta^{-1}(V).
\]
Then $C$ is a nonempty clopen of $S$ disjoint from $T$, and
$g_\beta:C\to V$ has the section $f_\alpha|_V$. In particular it is
surjective, contradicting \eqref{eq:topological-obstruction}.
\end{proof}

\begin{lemma}
\label{lem:freeness-factorization}
If $P\cong\Z[X]$ for a compact Hausdorff extremally disconnected
space $X$, then $\id_{\Z[T]}$ factors through $P^m$ for some finite
positive integer $m$.
\end{lemma}

\begin{proof}
Fix such an isomorphism and compose it with $P\hookrightarrow\Z[S]$.
By \cref{lem:finite-branches}, the resulting generator map has
finitely many continuous branches on clopen subspaces of $X$.
Let $H\subseteq S$ be the union of their images. It is closed,
being a finite union of compact subsets of a Hausdorff space.

For every $s\in S\setminus T$, the element
$[s]-[i(r(s))]\in P(*)$ is the image of a finite integral combination
of $X$-generators. Its $s$-coordinate is nonzero. Hence $s$ lies in
the image of one of the branches, and $S\setminus T\subseteq H$.
This subset is dense, so $H=S$.

Write $U_1,\ldots,U_m$ for the nonempty clopen branch domains, with
one copy for each branch. Their disjoint union
$Y=\coprod_{a=1}^m U_a$ maps continuously onto $S$. Here $m>0$ since
$S$ is nonempty. Extremal disconnectedness of $S$ gives a continuous
section, so $\Z[S]$ is a direct summand of $\Z[Y]$.
For each $a$, the clopen decomposition
$X=U_a\amalg(X\setminus U_a)$ makes $\Z[U_a]$ a direct summand of
$\Z[X]$. Consequently
\[
 \Z[Y]=\bigoplus_{a=1}^m\Z[U_a]
\]
is a direct summand of $\Z[X]^m\cong P^m$.
The original maps $r,i$ make $\Z[T]$ a direct summand of $\Z[S]$.
Composing these split diagrams gives the asserted factorization.
\end{proof}

\begin{proof}[Proof of \cref{prop:transfer}]
The group $P$ in \eqref{eq:projective-kernel} is nonzero compact
projective. If it were isomorphic to $\Z[X]$,
\cref{lem:freeness-factorization} would contradict
\cref{lem:no-scalar-factorization} with $n=1$.
\end{proof}

It remains to construct complete Boolean algebras satisfying
\cref{prop:transfer}. We first establish the tree-algebra obstruction
that will enforce its interval condition.

% End transfer.tex
% Begin trees.tex
\section{Tree representations and completion exclusion}\label{sec:trees}

We construct Boolean algebras from trees in such a way that omitting one tree
from a representation prevents its Boolean algebra from embedding into the
completion.  The relevant combinatorial input is Shelah's full super-bigness
theorem.  We state the part of that theorem used here, and then prove the
passage from its model chains to the required exclusion.

\subsection{The index trees and their Boolean algebras}

We use the standard trees of class $K^\omega_{\tr(3)}$ from
\cite[Definition~1.2]{Shelah2023}.  Here is their concrete form and language.
Fix a well-ordered set of ordinals.  A node is a sequence of length at most
$\omega$.  Every nonterminal coordinate is an increasing triple of ordinals;
the last coordinate of a nonempty finite node is a single ordinal.  The empty
sequence is a node.  If $\eta$ is a node and
\[
  \eta(i)=(\alpha_0,\alpha_1,\alpha_2),
  \qquad \alpha_0<\alpha_1<\alpha_2,
\]
then each sequence
\[
  (\eta\restr i)\mathbin{{}^\frown}\langle\alpha_m\rangle,
  \qquad m=0,1,2,
\]
is also a node.  The raw prefix $\eta\restr i$ need not itself be a node:
its last coordinate, if any, can still be a triple.

For $i<\omega$ and $m<3$, the unary function $\Res_i^m$ selects the displayed
node when $\eta(i)$ is a triple.  At the terminal scalar coordinate of a
finite node, and at every later level, it fixes the node.  It also fixes the
empty node.  This is the total restriction convention of
\cite[Definition~1.2(1)(B)($\theta$)]{Shelah2023}.
The remaining language consists of the root constant and the following
unary or binary relations:
\begin{itemize}
\item $P_i(\eta)$ means that $\eta$ has length $i$, for
      $i\leq\omega$;
\item $\eta\unlhd\nu$ means that $\eta$ is an initial segment of $\nu$;
\item $\Eq_i(\eta,\nu)$ means that
      $\eta\restr i=\nu\restr i$, for $i<\omega$; restriction of a
      sequence shorter than $i$ leaves that sequence unchanged;
\item $\eta<_1\nu$ means that both nodes have length $i+1$ for some
      $i<\omega$, have the same prefix of length $i$, and their last scalar
      coordinates are in increasing order;
\item $\operatorname{Suc}_{i,m}(\eta,\nu)$ means that $\nu(i)$ is a
      triple and $\eta=\Res_i^m(\nu)$.
\end{itemize}
In particular, every sibling fiber is well ordered.  Throughout, the
quantifier-free type of a tuple of indices is taken in this full language,
including formulas with restriction terms.
In this concrete syntax, the literal initial-segment relation satisfies
\[
 v\unlhd w\quad\Longleftrightarrow\quad v=\varnothing\ \text{or}\ v=w.
\]
Indeed, a nonempty finite node ends in a scalar, whereas a longer node
has a triple at that position; two infinite nodes can be initial
segments of one another only if they are equal.

For such a tree $I$, let $\BA(I)$ be the Boolean algebra generated by
$x_v$, $v\in I$, freely subject to
\begin{equation}\label{eq:tree-relations}
 x_\eta\leq x_{\Res_n^0(\eta)},
 \qquad
 x_\eta\wedge x_{\Res_n^1(\eta)}\wedge x_{\Res_n^2(\eta)}=0
 \quad(\eta\in P_\omega(I),\ n<\omega).
\end{equation}
This is \cite[Definition~2.1(3)]{Shelah2026}.  Each generator is nonzero.
The following direct description explains precisely how the tree language
controls finite Boolean relations.

\begin{lemma}\label{lem:tree-relations}
The Boolean quantifier-free type of a finite tuple of generators of
$\BA(I)$ is determined by the joint quantifier-free type of their indices.
\end{lemma}

\begin{proof}
For a finite tuple of Boolean elements, a \emph{signed cell} is the meet
obtained by choosing each element or its complement.  Every Boolean equation
in the tuple is determined by which signed cells vanish.  It therefore
suffices to determine when a signed cell in the generators is nonzero.
Declare its positive
generators true, and also declare every zeroth restriction of each positive
infinite generator true.  The resulting requirements are consistent exactly
when both of the following conditions hold:
\begin{enumerate}
\item no generator required to be false has been declared true;
\item for no positive infinite generator and no level have both its first
      and second restrictions been declared true.
\end{enumerate}
Necessity follows from \eqref{eq:tree-relations}.  For sufficiency, assign
false to every generator not yet declared true.  All the defining relations
hold: relations whose infinite generator is false impose no condition, and
those whose infinite generator is true are precisely the ones already
checked.  This gives a homomorphism to the two-element Boolean algebra under
which the cell is true.

The two consistency tests use only equalities among the given indices and
their restriction terms.  Thus they are preserved by equality of the full
joint index types and determine the Boolean quantifier-free type.  The same
tests show that every generator is nonzero and that distinct generators
can be separated: when one of the two is finite, declare that generator
true and the other false; when both are infinite, declare one true and
the other false.  In each case the tests are satisfied.
\end{proof}

\subsection{Ordinary representations and the combinatorial input}

A code is a term whose parameters are indices from a tree $J$.  Two tuples
of codes are \emph{similar} if their term shapes agree and their parameter
tuples have the same joint quantifier-free type in $J$.  Repeated parameter
positions are retained, so the joint type includes all equalities between
positions in different codes.

An \emph{ordinary representation} of a Boolean algebra $F$ over $J$ is a
map $\rho$ assigning a code to each element of $F$, such that similar tuples
of codes come from tuples with the same Boolean quantifier-free type.  This
is the representation condition of \cite[Definition~2.1(c),(e),(g)]{ShelahE59}.
It is a condition on the information encoded by terms, rather than a
homomorphism into a term algebra.  The information consists of term shapes
and index types; it contains no additional ordering of terms.

For finitary representations we fix a countable finitary vocabulary, so
each code has finitely many subterms and index parameters.  We reserve
disjoint families of symbols for initial constants, relative interval
codes, finite Boolean expressions, and local copies of tree generators.
There are countably many symbols of each finite arity available for these
purposes.  Completion will use a separate countably ary symbol.  All the
countable codes thereby obtained belong to
$M_{\cc,\aleph_1}(J)$, the free term algebra with at most $\cc$ symbols of
each countable arity.  Here and below
\[
 \cc=2^{\aleph_0},\qquad \lambda=(2^\cc)^+.
\]

For a family of standard trees, its specialized sum identifies the empty
roots and places first coordinates from different trees in disjoint,
ordered ordinal blocks.  More precisely, after choosing a common ordinal
bound $\zeta_*$ for their coordinates, the first coordinate in the
$\gamma$th summand is shifted into the block
$[\zeta_*\gamma,\zeta_*(\gamma+1))$; the shift is applied componentwise
when that coordinate is a triple.  The other coordinates are unchanged.
This is the sum of \cite[Definition~1.4(4)]{Shelah2023}.

\begin{theorem}[Shelah's full super-bigness input]\label{thm:tree-input}
There is a family $(I_\gamma)_{\gamma<\lambda}$ of standard
$\tr(3)$ trees, each of cardinality $\lambda$, such that
$I_\gamma$ is $(\cc,\cc)$-super-unembeddable into
\[
  J_\gamma=\sum_{\substack{\delta<\lambda\\\delta\ne\gamma}}I_\delta
\]
with the default equivalence relation on successor positions.  Moreover,
$\BA(I_\gamma)$ has cardinality $\lambda$ and is Knaster.
\end{theorem}

The existence assertion is
\cite[Definitions~1.4(1),(3),(4), 1.6(2), and Lemma~1.11(1)]{Shelah2023}.
Its cardinal hypothesis holds because $\lambda$ is regular and
$\lambda>\cc^\cc=2^\cc$.  The Knaster assertion is
\cite[Claim~2.6(3)]{Shelah2026}, applied with height $\omega$ and chain
condition cardinal $\aleph_1$.  The cardinality follows from the
presentation and the distinctness of generators proved in
\cref{lem:tree-relations}.  Fix this family and put
\[
 E_\gamma=\BA(I_\gamma)\qquad(\gamma<\lambda).
\]
Recall that Knaster means that every
uncountable family of nonzero elements has an uncountable pairwise
compatible subfamily; in particular these algebras are ccc.

We spell out the model-chain conclusion of super-unembeddability needed
below.  Suppose that $I$ is $(\cc,\cc)$-super-unembeddable into $J$ as in
the theorem.  Fix a sufficiently large regular cardinal $\chi$.
Write $H(\chi)$ for the sets of hereditary cardinality less than $\chi$,
and fix a well-order $<^*_\chi$ of $H(\chi)$.  Given any prescribed
parameter $x\in H(\chi)$ and a map
\[
 f_1:I\longrightarrow {}^{<\cc}J,
\]
there are elementary submodels
\begin{equation}\label{eq:model-chain}
 M_n\prec N_n\prec M_{n+1}\prec(H(\chi),\in,<^*_\chi)
 \qquad(n<\omega)
\end{equation}
and $\eta\in P_\omega(I)$ with the following properties.  We include
$I,J,f_1$ among the prescribed parameters.  The prescribed
objects belong to $M_0$, and $\cc\subseteq M_0$.  For all sufficiently
large $n$, the three nodes $\Res_n^m(\eta)$ lie in $N_n\setminus M_n$,
and the three sequences
\begin{equation}\label{eq:three-sequence-types}
 f_1(\Res_n^m(\eta)),\qquad m=0,1,2,
\end{equation}
have the same length and the same atomic type over $J\cap M_n$.
Equality of atomic types here includes the truth values of all atomic
formulas, and therefore determines the quantifier-free type.

There is one further condition.  Set $U=\bigcup_n M_n$.  For each
$w\in P_\omega(J)$, all restrictions $\Res_i^m(w)$ which belong to $U$
belong together to some $M_r$.  This follows from clause~(vi) of
\cite[Definition~1.4(1)]{Shelah2023}.  In the printed clause the union of
the finite restriction nodes is interpreted as a union of sets.  If a
finite restriction node lies in $U$, all its finitely many members lie in
$U$ and hence in the one model bounding that union.  Closure under finite
set formation puts the restriction node itself in the same model.

The default equivalence relation in
\cite[Definition~1.4(3)]{Shelah2023} compares every pair of singleton
successor positions.  It is this part of the hypothesis that gives all
three sequences in \eqref{eq:three-sequence-types} the same type.

For subsequent constructions choose a marker
$t_\gamma\in P_1(I_\gamma)$ in each tree.  In a specialized sum,
\begin{equation}\label{eq:marker-order}
 t_\gamma<_1t_\delta\quad\Longleftrightarrow\quad\gamma<\delta.
\end{equation}
Thus equality and chronology of markers are visible in the ordinary index
type.  Relabelling any subfamily into the sum omitting one unused tree
preserves all joint types, because it preserves the order of the blocks
and the structure within each block.

\begin{proposition}[Completion exclusion]\label{prop:exclusion}
Suppose $I$ is $(\cc,\cc)$-super-unembeddable into a standard tree $J$
with the default equivalence relation.  Let $F$ be a dense subalgebra of a
Boolean algebra $C$.  If $F$ has an ordinary finitary representation over
$J$ extending to an ordinary representation of $C$ by countable terms,
then $\BA(I)$ does not embed into $C$.
\end{proposition}

The obstruction is the pair of relations in \eqref{eq:tree-relations}.
We will choose a positive element of $F$ below the image of one infinite
generator.  To transfer its inequality with the zeroth successor to both
other successors, we need similarity after adjoining the finite parameters
of this chosen element.  The next lemma supplies that parameter transfer;
the following lemma applies it simultaneously to all three successors.

\subsection{Adjoining finitely many external indices}

\begin{lemma}\label{lem:external-parameters}
Let $J$ be a standard $\tr(3)$ tree, and let $(M_n,N_n)_{n<\omega}$ be
a chain as in \eqref{eq:model-chain}, with $J\in M_0$ and with the
restriction-boundedness condition just stated.  Put $U=\bigcup_nM_n$.
For every finite $W\subseteq J$ there exists $n_0<\omega$ such that,
for $n\geq n_0$, any two equally indexed tuples from $J\cap U$ having
the same quantifier-free type over $J\cap M_n$ still have the same joint
type after adjoining the fixed tuple enumerating $W$.
\end{lemma}

\begin{proof}
Every function in the tree language is unary, and every relation has arity
at most two.  Moreover, an iterated restriction of a node is either the
node itself or a single restriction: restricting at an earlier level
selects the earlier restriction, whereas restricting at or beyond a
terminal level fixes that node.  Initial-segment tests reduce to root
tests and equality by the displayed description above. It therefore
suffices to handle the other mixed
atomic relations between an internal node, meaning a node of $J\cap U$,
and a node or restriction arising from one fixed $w\in W$.  Relations
involving only one of the two tuples already agree.  We show that all
mixed relations are decided by finitely many fixed internal nodes,
except in one case where the model-chain condition puts all necessary
nodes in a common model.

\medskip\noindent
\emph{An internal node $w$.}
Choose $r$ with $w\in M_r$.  Since the restrictions are definable from
the tree and their finite indices, all restrictions of $w$ also belong
to $M_r$.  The mixed relations are then already among the tests over
$J\cap M_n$ for every $n\geq r$.

\medskip\noindent
\emph{An external node whose restrictions are all internal.}
Suppose $w\notin U$, but every $\Res_i^m(w)$ belongs to $U$.
Then $w$ is infinite, since a finite node is its own terminal restriction.
The restriction-boundedness condition supplies one $M_r$ containing all
these restrictions.  Equality of $w$ with an internal node is false.
Every finite prefix test against $w$ is decided by its restrictions; for
example,
\[
 \Eq_i(z,w)\quad\Longleftrightarrow\quad
 \Eq_i(z,\Res_i^0(w)).
\]
Length predicates take care of nodes shorter than the prefix being tested.
A successor relation is an equality to the appropriate restriction with
the prescribed levels.  Sibling comparisons with $w$ are false because
$w$ is infinite.  These statements also handle
every restriction term applied to $w$, since the restrictions themselves
belong to $M_r$.

\medskip\noindent
\emph{An external node with a first missing restriction.}
Let $k$ be the least level at which some
\[
 r_m=\Res_k^m(w)
\]
does not belong to $U$.  The earlier restrictions determine the raw prefix
$\sigma=w\restr k$, so $\sigma\in U$; for $k=0$ this is the empty
sequence.  The sibling fiber above $\sigma$ is nonempty.  By elementarity
there is an actual node $q_\sigma\in J\cap U$ in this fiber, of length
$k+1$.  Fix such a node.  The formula $\Eq_k(-,q_\sigma)$ detects the
prefix $\sigma$, with length predicates distinguishing shorter nodes.
The parameter used here is the index node $q_\sigma$, not the possibly
non-index sequence $\sigma$.

Choose a model containing $q_\sigma$, all earlier restrictions of $w$,
and all of the $r_m$ which belong to $U$.  For each remaining $r_m$,
consider the internal nodes in its sibling fiber strictly above it.
If this set is nonempty, its well ordering gives a least element $d_m$;
enlarge the model to contain $d_m$.  If the set is empty, no parameter is
needed.  There are at most three choices to make.  For an internal node
$z$ in that fiber, equality $z=r_m$ is false, and
\begin{equation}\label{eq:first-exit-cut}
 z<_1r_m\quad\Longleftrightarrow\quad z<_1d_m
\end{equation}
when $d_m$ exists.  Indeed, $d_m$ is the first internal node strictly
above $r_m$, and $r_m$ itself is not internal.  If no $d_m$ exists, every
internal node in the fiber is below $r_m$.  The reverse inequalities are
determined by the same data.  Thus the mixed tests in the first missing
fiber reduce to tests over the chosen model.

Suppose next that $w$ has a coordinate after this fiber.  No internal node
can have raw prefix $w\restr(k+1)$: extracting its finitely many
coordinates would recover the missing restriction $r_m$ inside $U$.
Consequently every genuinely deeper restriction of $w$ is external.
Equality with an internal node is false, prefix equality beyond the first
exit is false, and sibling comparisons in a deeper fiber are false.
Prefix comparisons up to the exit are determined by $q_\sigma$ and the
earlier restrictions.  Successor relations in either
orientation are again equalities with restrictions, together with level
conditions.

If instead $w$ is already terminal at the first missing level, its later
restrictions repeat $w$.  All comparisons with these repeated terminal
restrictions are governed by the same cut \eqref{eq:first-exit-cut}; they
are not comparisons in deeper fibers.  This also covers repeated
terminal restrictions after any subsequent application of a restriction
function.

We have accounted for every relation and every restriction term in the
language.  For each $w\in W$ the required parameters lie in one model;
take a common model for the finitely many $w$.  The descriptions apply
to every internal node, independently of the length of the variable
tuple.  Each atomic formula uses finitely many entries, so equality of
the original types preserves all the mixed atomic formulas.  It then
preserves their Boolean combinations, as required.
\end{proof}

\begin{lemma}[Simultaneous successor similarity]\label{lem:three-successors}
Let $I$ be $(\cc,\cc)$-super-unembeddable into a standard tree $J$
with the default equivalence relation.  Suppose
\[
 f:I\longrightarrow M_{\cc,\aleph_1}(J)
\]
assigns countable terms, and assigns finite terms with finite index support
to the nodes of $P_\omega(I)$.  Then there are $\eta\in P_\omega(I)$
and $n<\omega$ such that the three pairs
\[
 \bigl(f(\eta),f(\Res_n^m(\eta))\bigr),\qquad m=0,1,2,
\]
are similar.
\end{lemma}

\begin{proof}
There are at most $\cc$ countable term shapes: each syntax tree is
countable, the vocabulary has size at most $\cc$, and
$\cc^{\aleph_0}=\cc$.  Enumerate the shapes by ordinals $\xi<\cc$.
For each term, put its parameters in the first $\omega$ positions of a
sequence, padding unused positions with the root of $J$.  The shape
records which positions are used.  For a term of shape $\xi$, continue
the padding to length
\[
 L_\xi=\omega\cdot(\xi+1)<\cc.
\]
These lengths are distinct.  The inequality holds because each such
ordinal has cardinality less than $\cc$; it does not require regularity
of $\cc$.  We have obtained a map
$f_1:I\to{}^{<\cc}J$ that records both shapes and parameters.

Apply the model-chain conclusion to $f_1$, including the map and the
coding choices in the prescribed parameter.  Let $\eta$ be its witness
and let $W$ be the finite support of $f(\eta)$.  Choose one $n$ large
enough for \cref{lem:external-parameters} and for
\eqref{eq:three-sequence-types}.  The three successor sequences have
equal lengths and equal types over $J\cap M_n$.
Each sequence belongs to $N_n$, since its argument and $f_1$ do.
Every one of its coordinates also belongs to $N_n$: its length is below
$\cc$ and every ordinal below $\cc$ belongs to $M_0$.
Thus these are tuples of indices in $J\cap U$.

By \cref{lem:external-parameters}, their types remain equal after
adjoining $W$.  Equality of lengths gives equality of their term shapes.
Together with the fixed shape and support of $f(\eta)$, this proves
similarity of all three pairs at this same $\eta$ and $n$.
\end{proof}

\begin{proof}[Proof of \cref{prop:exclusion}]
Suppose $h:\BA(I)\hookrightarrow C$ is an embedding.  For each finite
node $v$, let $f(v)$ be the countable code of $h(x_v)$.  For each
$\eta\in P_\omega(I)$, density and $h(x_\eta)>0$ give an element
\[
 0<a_\eta\in F,\qquad a_\eta\leq h(x_\eta).
\]
Use the retained finite code of $a_\eta$ for $f(\eta)$.
Apply \cref{lem:three-successors} and write
$\nu_m=\Res_n^m(\eta)$ for the resulting successors.
The defining relations give
\[
 a_\eta\leq h(x_\eta)\leq h(x_{\nu_0}).
\]
Similarity and ordinary representability transfer the last inequality
with $a_\eta$ to both other successors.  Hence
\[
 0<a_\eta
 \leq h(x_\eta)\wedge h(x_{\nu_1})\wedge h(x_{\nu_2})=0,
\]
contrary to \eqref{eq:tree-relations}.
\end{proof}

\section{Represented intervals and local insertions}
\label{sec:completion-codes}

The construction will repeatedly use intervals whose units belong to a
completed algebra but not necessarily to its represented dense subalgebra.
We first produce finite codes for a dense algebra in each such interval.
We then give countable codes for a completion.  Keeping these two roles
separate allows later constructions to retain finite codes while
\cref{prop:exclusion} is applied to the completed algebra.

\subsection{Dense representations and completion codes}

\begin{lemma}\label{lem:interval-representation}
Let $F$ be a dense subalgebra of a complete Boolean algebra $C$, with an
ordinary finitary representation $\rho$ over $J$.  For every $0<c\in C$,
the interval $C\restr c$ has a dense subalgebra with an ordinary finitary
representation over the same $J$.
\end{lemma}

\begin{proof}
If $c\in F$, take $F\restr c$.  With a reserved binary symbol $R$,
assign $a\in F\restr c$ the code
\[
 R\bigl(\rho(c),\rho(a)\bigr).
\]
Relative complementation is $a\mapsto c\wedge\neg a$, so every
quantifier-free test in the interval is an old Boolean test with the
additional entry $c$.  The displayed code contains exactly the data
needed for these tests.

Now suppose $c\notin F$.  Put
\[
 H=\{a\in F:a\leq c\},\qquad
 G=H\cup\{c-a:a\in H\},
\]
where $c-a=c\wedge\neg a$.  The set $H$ is an ideal of the Boolean
ring $F$, and $G$ is its unitization with unit $c$.  In particular, $G$
is a Boolean subalgebra of $C\restr c$ and is dense there because it
contains $H$.  Its two displayed forms are disjoint: if $a=c-b$ with
$a,b\in H$, then $c=a\vee b\in F$, a contradiction.

Use reserved constants for $0$ and $c$.  For the other elements, use
distinct unary symbols $U_+$ and $U_-$ and assign
\[
 a\longmapsto U_+(\rho(a)),\qquad
 c-a\longmapsto U_-(\rho(a)).
\]
To verify the representation condition, reduce a Boolean equation to the
vanishing of signed cells.  After interpreting the tags, such a cell is
an intersection of elements of $H$ and factors of the form $c-b$ with
$b\in H$.  If at least one $H$-factor occurs, its vanishing is an old
Boolean equation in $F$, since that factor already lies below $c$.
If none occurs, the cell is
\[
 c-\bigvee_{j<r}b_j,
\]
which is nonzero: its vanishing would put $c=\bigvee_{j<r}b_j$ in $F$.
Constants $0$ and $c$ are dealt with directly.  The old ordinary
representation therefore determines every required test.

In both cases the outer symbol identifies the code of an entire element
of the relative algebra.  In particular, the old subterms inside a
relative code are used to evaluate old Boolean equations; they need not
separately code elements of the relative algebra.
\end{proof}

\begin{lemma}\label{lem:completion-representation}
Let $F$ be a ccc Boolean algebra with an ordinary finitary representation
$\rho$ over $J$, and let $C$ be its completion.  There is an ordinary
representation
\[
 \widehat\rho:C\longrightarrow M_{\cc,\aleph_1}(J)
\]
by countable terms which retains the given code of every element of $F$.
\end{lemma}

\begin{proof}
For each $x\in C$, choose countable families in $F$ such that
\begin{equation}\label{eq:two-sided-approximants}
 x=\bigvee_{n<\omega}a_{x,n},\qquad
 \neg x=\bigvee_{n<\omega}b_{x,n}.
\end{equation}
To obtain the first family, choose a maximal pairwise disjoint family of
nonzero elements of $F$ below $x$.  It is countable by ccc, and its join
is $x$ by density and maximality.  Do the same below $\neg x$, and pad
finite families by zeros.

For $x\in F$ set $\widehat\rho(x)=\rho(x)$.  For $x\notin F$, use a
reserved countably ary symbol $\mathsf H$ and set
\[
 \widehat\rho(x)
 =\mathsf H\bigl(\rho(a_{x,0}),\rho(b_{x,0}),
                 \rho(a_{x,1}),\rho(b_{x,1}),\ldots\bigr).
\]
The symbol $\mathsf H$ distinguishes these codes from retained finite
codes.  Each has countably many finite constituent terms, and hence only
countably many subterms and index parameters.

Consider a signed cell in a finite tuple of elements of $C$.  For an
entry outside $F$, use the positive or negative approximating join in
\eqref{eq:two-sided-approximants}, according to its sign.  Keep entries
from $F$ as old Boolean expressions.  Finite meets distribute over
arbitrary joins in the complete Boolean algebra $C$.  Thus the cell is
nonzero exactly when some selection of one approximating piece for each
new entry has nonzero meet with the retained old expressions.  Each
selection is a quantifier-free test on a finite tuple from $F$, and its
codes are visible subterms of the displayed completion codes.

Equal term shapes and equal joint index types preserve the data for
every one of these finite tests.  Ordinary representability of $F$
therefore determines the vanishing of the original cell and hence every
Boolean equation.  This proves the representation condition.  It also
proves the condition for countable tuples, since each quantifier-free
formula uses only finitely many entries.
\end{proof}

In subsequent constructions the finite representation of a dense algebra
is retained.  The countable extension in
\cref{lem:completion-representation} is used when applying
\cref{prop:exclusion} to its completion.

% End trees.tex
% Begin surgery.tex
\subsection{Local insertions and complementary intervals}\label{sec:surgery}

The construction will adjoin countably many free copies of a fresh tree
algebra below an element of a completed Boolean algebra.  That element
need not belong to the dense
subalgebra carrying the finite representation.  The following lemma realizes
the extension using countably many elements of the dense subalgebra.  Its
finite product decomposition supplies the representation at successor stages;
regular unions supply the limit stages.

For each $\gamma<\lambda$, put
\[
 G_\gamma=*_{m<\omega}E_\gamma^{(m)},
\]
where $E_\gamma^{(m)}$ is a copy of $E_\gamma$ and $*$ denotes the free
product of unital Boolean algebras.  Write $x_\eta^{(m)}$ for the generator
indexed by $\eta\in I_\gamma$ in the $m$th copy.

\begin{lemma}\label{lem:countable-free-product}
The algebra $G_\gamma$ has cardinality $\lambda$ and is Knaster.
\end{lemma}

\begin{proof}
It contains a copy of $E_\gamma$, and its elements are finite Boolean
expressions in countably many copies of a set of cardinality $\lambda$.
Hence $|G_\gamma|=\lambda$.

Every positive element contains a positive rectangle
$\bigwedge_{m\in L}e_m$, where $L\subseteq\omega$ is finite and
$0<e_m\in E_\gamma^{(m)}$.  Given an uncountable family of positive
elements, choose such rectangles and thin to an uncountable subfamily
with the same $L$.  For each of the finitely many $m\in L$, use the
Knaster property of $E_\gamma^{(m)}$ to thin again so that the $m$th
factors are pairwise compatible.  Two resulting rectangles have nonzero
meet because their within-copy meets are nonzero and the different
copies are freely independent.  The corresponding original elements
are therefore pairwise compatible.
\end{proof}

Recall that an inclusion of Boolean algebras is \emph{regular} if every
maximal antichain of the smaller algebra remains maximal in the larger one.
A regular embedding between complete Boolean algebras is complete.  Indeed,
an arbitrary join can be computed using a maximal disjoint refinement of the
elements being joined, together with the complement of their join.

\begin{lemma}\label{lem:surgery}
Let $C$ be a complete ccc Boolean algebra, and let $F\subseteq C$ be a dense
subalgebra with an ordinary finitary representation.  Suppose that its codes
use only trees $I_\delta$ with $\delta<\gamma$, and let $0<u\in C$.  Put
\begin{equation}\label{eq:completed-surgery}
 C'=(C\restr\neg u)\times
       \widehat{(C\restr u)*G_\gamma}.
\end{equation}
Here $*$ denotes the free product of unital Boolean algebras and the hat
denotes Boolean completion.  Then $C'$ is ccc and the canonical inclusion
$C\hookrightarrow C'$ is complete.

For the representation assertion, assume the following history of the
codes.  The initial dense algebra has an ordinary representation whose
\emph{whole codes} have distinguished outer symbols: either the nullary
symbols of the initial algebra, or the interval and unitization symbols of
Lemma~\ref{lem:interval-representation}.  The remaining codes have been
formed by the finite normal-form prescription in the proof below, using
strictly increasing tree indices and a disjoint collection of outer symbols.
Every such stage index is greater than all tree indices in the initial codes.
Decoding stops at a whole initial code and uses its given ordinary
representation.  In particular, a subterm inside an initial interval code is
not decoded as a later stage of this history.

Under this hypothesis, $C'$ has a dense subalgebra $F'\supseteq F$ such that
$F\subseteq F'$ is regular and the given representation extends to an
ordinary finitary representation of $F'$.  The extension uses only the old
trees and $I_\gamma$, retains all old codes, and has the same code-history
property.  These representations remain coherent under iteration and
continuous unions of the dense algebras.
\end{lemma}

\begin{proof}
We first identify a dense algebra through its finite product decompositions.
We then encode these decompositions and show that the tree types of the
codes determine every Boolean equation.  The latter assertion is proved by
successively removing the greatest construction index present in a finite
tuple of codes.

\paragraph{The complete extension.}
Every positive element of $(C\restr u)*G_\gamma$ contains a rectangle
$d\wedge e$ with $0<d\leq u$ in $C$ and $0<e\in G_\gamma$.
An old maximal antichain meets $d$, and hence meets this rectangle.  Together
with the unchanged factor below $\neg u$, this shows that the ordinary
surgery is a regular extension of $C$.  Passing to its completion preserves
regularity, so the inclusion in \eqref{eq:completed-surgery} is complete.
The same rectangle argument gives the regularity assertion in
\cite[Claim~2.9(2)]{Shelah2026}.

The algebra $G_\gamma$ is Knaster by
Lemma~\ref{lem:countable-free-product}.  In an uncountable family of
positive rectangles, therefore, uncountably many fresh factors are
pairwise compatible.  Since
$C\restr u$ is ccc, two of the corresponding old factors are compatible.
Their rectangles are compatible.  Thus the free product is ccc, as are its
completion and the finite product $C'$; compare
\cite[Claim~2.9(1)]{Shelah2026}.

\paragraph{Finite products below disjoint old elements.}
Choose a maximal disjoint family $(a_n)_{n\in N}$ of positive elements of
$F$ below $u$.  Here $N\subseteq\omega$ may be finite.  Density of $F$ and
the ccc imply
\[
 \bigvee_{n\in N}^{C}a_n=u.
\]
In $C'$, let $F'$ be generated by $F$ and all the elements $a_n\wedge e$,
where $n\in N$ and $e\in G_\gamma$.  We regard the fresh algebra as lying
in the second factor of \eqref{eq:completed-surgery}, with relative unit
$u$.

For a finite subset $L\subseteq N$, set $v=\bigvee_{n\in L}a_n$, and let
$F_L$ be the subalgebra generated using only the local copies indexed by
$L$.  Restriction to these finitely many disjoint elements gives
\begin{equation}\label{eq:finite-band-product}
 F_L\cong
 (F\restr\neg v)\times
 \prod_{n\in L}\bigl((F\restr a_n)*G_\gamma\bigr).
\end{equation}
On the $n$th factor, old and fresh nonzero elements have nonzero meet by
the defining independence of the free product in
\eqref{eq:completed-surgery}.  This identifies that factor with
$(F\restr a_n)*G_\gamma$.  Conversely, a tuple in the displayed finite
product can be pasted using the elements $a_n$ and $\neg v$ of $F$.
This proves both surjectivity and injectivity of the displayed map.
In particular, the use of one copy of $G_\gamma$ in the completed
extension introduces no further relations between the different factors of
\eqref{eq:finite-band-product}.

The algebra $F'$ is dense in $C'$.  For a positive rectangle $d\wedge e$
in the second factor, some $a_n\wedge d$ is positive.  Choose
$0<a\in F$ with $a\leq a_n\wedge d$.  Then
\[
 0<a\wedge e=a\wedge(a_n\wedge e)\in F'
\]
lies below the rectangle.  Density in the first factor follows from
density of $F$ in $C$.  Every positive element of the completion contains
a positive element of the uncompleted free product, so these observations
prove density in all of $C'$.
Finally, an antichain maximal in $F$ has join $1$ in $C$.  Completeness of
$C\hookrightarrow C'$ implies that it still has join $1$ in $C'$, and
hence is maximal in $F'$.  Thus $F\subseteq F'$ is regular.

\paragraph{Finite normal forms and their codes.}
Fix countably many symbols for finite Boolean normal forms and binary
symbols $T_{n,m}$, $n,m<\omega$, disjoint from the outer symbols of the initial
codes.  Use the first-level marker $t_\gamma\in I_\gamma$ fixed in
Section~\ref{sec:trees}.  The occurrence of the generator
$x_\eta^{(m)}$ in the $n$th local copy is recorded by the subcode
\begin{equation}\label{eq:local-tree-code}
 T_{n,m}(t_\gamma,\eta).
\end{equation}
This is not by itself the code of the element $a_n\wedge x_\eta^{(m)}$:
the old support unit must also be recorded.  With a reserved ternary
symbol $\mathsf{Gen}$, a full code for that element is
\[
 \mathsf{Gen}\bigl(t_\gamma,\rho(a_n),
                         T_{n,m}(t_\gamma,\eta)\bigr).
\]
For each new element of $F'$, choose a normal form in some $F_L$ as in
\eqref{eq:finite-band-product}.  Its code records the finite set of local
indices, the Boolean expressions in the fresh generators, the retained
codes of all old coefficients, and the retained codes of the units $a_n$
used in the expression.  The marker $t_\gamma$ occurs in a designated
position in the outer code, as well as in
\eqref{eq:local-tree-code}.  There are only countably many possible finite
normal-form symbols.  The tree parameter $t_\gamma$ distinguishes the
stage, so this vocabulary can be reused at all stages.  Elements already
in $F$ retain their old codes.

We verify that this is an ordinary representation, including for tuples
whose elements were introduced at different stages.  Suppose two finite
tuples of codes have the same term shapes and the same joint
quantifier-free type of their tree parameters.  Starting from the outer
codes, stop whenever an initial whole code is reached.  Outside those
initial codes, only finitely many construction stages occur.  Equality of
the designated markers groups the occurrences belonging to one stage.
The relation $<_1$ between the markers orders these groups by construction
time, by the ordering convention of Section~\ref{sec:trees}.  The two
tuples therefore have corresponding greatest stages.

To decide a Boolean equation in either tuple, expand its difference into
signed cells.  At the greatest stage, refine each cell by all the finitely
many participating units $a_n$ and their common complement.  Formula
\eqref{eq:finite-band-product} describes the resulting restrictions.
On its $n$th factor, a Boolean expression is a finite union of rectangles
whose old factors are Boolean expressions in the recorded old
coefficients and units, and whose fresh factors are Boolean expressions
in the recorded generators of $G_\gamma$.  Such a rectangle is zero
exactly when one of its factors is zero.  On the complementary factor,
only the old expressions remain.

For a fresh-factor test, expand its expression into signed cells and
group the generators in each cell by their copy index $m$.  Freeness of
the copies implies that a cell is nonzero exactly when every within-copy
cell is nonzero.  Lemma~\ref{lem:tree-relations} determines each
within-copy test from the joint tree type.  The symbols $T_{n,m}$
identify both the old support band and the free copy on the two sides.
All remaining tests are Boolean equations in old coefficients and old
units, whose codes are recorded subterms.  Thus
corresponding tests reduce to corresponding tests with the greatest stage
removed.  Repeating this reduction terminates, because a finite term has
only finitely many stage occurrences.  At termination the ordinary
representation of the initial dense algebra decides the remaining tests.
Consequently the original Boolean equations, and their negations, have
the same truth value for the two tuples.  This is the required ordinary
representation property.

For a relative initial algebra, this last step applies to its whole
interval or unitization codes.  Lemma~\ref{lem:interval-representation}
evaluates their Boolean equations using the original old representation.
The internal coefficients of such a code serve that evaluation; they are
not asserted to be elements of the relative initial algebra.  The
distinguished outer symbols specify exactly where the finite-stage
reduction stops.

\paragraph{Continuous unions of the dense algebras.}
At a limit stage take the union of the preceding dense algebras, with their
retained codes, and then take its Boolean completion.  Any two finite
tuples under comparison lie together in one earlier dense algebra, so
the ordinary representation property passes to the union.  The dense
algebras form a continuous regular chain.  Its union is ccc by
\cite[Claim~2.10]{Shelah2026}, and every earlier dense algebra is regular
in that union.  The induced maps between the completions are therefore
complete embeddings.  This proves the asserted coherence at successor
and limit stages.
\end{proof}

The dense algebras in this construction form the continuous chain; at a
limit the completed algebra is the completion of their union.  Countable
representations of these completions are supplied separately by
Lemma~\ref{lem:completion-representation}.  Later applications of
Lemma~\ref{lem:surgery} retain the finite dense codes.  In a relative
iteration, a stage whose cut unit is zero leaves both the algebra and
its codes unchanged and skips the corresponding tree index.

\begin{lemma}\label{lem:complementary-omission}
Consider an iteration as in Lemma~\ref{lem:surgery}, using fresh tree
indices in strictly increasing order and the first-level markers of
Section~\ref{sec:trees}.  Suppose that $G_\gamma$ is adjoined at the old
completed unit $u_\gamma$.  If $B$ is the final completion, then there is
no unital Boolean embedding
\[
 E_\gamma\hookrightarrow B\restr\neg u_\gamma,
\]
where the target has relative unit $\neg u_\gamma$.
\end{lemma}

\begin{proof}
We construct a dense ordinary finitary representation of the complementary
interval whose parameters lie in the omitted-tree sum $J_\gamma$.
Its extension to the completion will then satisfy the hypotheses of
Proposition~\ref{prop:exclusion}.

The assertion is immediate if $\neg u_\gamma=0$.  Otherwise write
$c=\neg u_\gamma>0$ in the completed algebra immediately before the
$\gamma$th surgery.  Lemma~\ref{lem:interval-representation} gives a
dense subalgebra of that interval with an ordinary finitary representation
using only trees of index less than $\gamma$.  Take its whole interval
codes as the initial codes for a new relative construction.  The surgery
at $u_\gamma$ leaves the interval below $c$ unchanged.

Suppose a later global surgery adjoins $G_\delta$ below $u_\delta$ in an
old complete algebra $C$.  Restricting its result to $c$ gives
\begin{equation}\label{eq:relative-completed-surgery}
 (C\restr(c\wedge\neg u_\delta))\times
 \widehat{(C\restr(c\wedge u_\delta))*G_\delta}.
\end{equation}
For a zero element $c\wedge u_\delta$, this means that the interval is
unchanged.  For a nonzero element, the identity follows by splitting the
old factor of the free product at its clopen $c$.  Under Stone duality a
free product becomes a product of compact spaces, and the clopen
partition of the old space gives a finite disjoint decomposition of that
product.  Boolean completion commutes with this finite decomposition.
Thus \eqref{eq:relative-completed-surgery} is precisely the completed
surgery in the old relative algebra $C\restr c$, with cut unit
$u_\delta\wedge c$.

Apply Lemma~\ref{lem:surgery} at each of these relative successor stages,
choosing a countable disjoint family from the current relative dense
algebra below the relative cut unit.  This constructs a separate chain of
dense algebras in the intervals below $c$.  Its initial whole codes are
those of Lemma~\ref{lem:interval-representation}, and every later tree
index is greater than every index in those codes.  Hence the code-history
hypothesis of Lemma~\ref{lem:surgery} holds throughout this relative
construction.

At a limit stage, its relative union is dense in the corresponding global
interval.  Indeed, if $0<b\leq c$ lies in the global completion at that
stage, the global dense union contains $0<a\leq b$.  Both $a$ and $c$
belong to an earlier completed stage, after the relative construction
has begun.  The relative dense algebra at that earlier stage contains
$0<g\leq a$.  This element $g$ belongs to the relative union and lies
below $b$.  The complete embeddings therefore identify the completion of
the relative union with the global interval.  The same argument applies
to the final completion $B$.

We have obtained a dense ordinarily finitarily represented subalgebra of
$B\restr c$ whose parameters belong to the other tree components and
the common root.  In particular, no nonroot parameter from the
$I_\gamma$ component, including its marker $t_\gamma$, occurs.
The initial relative construction precedes $I_\gamma$, and each later
relative surgery introduces only its own tree and marker; its countably
many free copies are distinguished by vocabulary symbols, not by
additional trees.
Relabel the allowed trees into the full omitted-tree sum $J_\gamma$ of
Section~\ref{sec:trees}; this preserves joint tree types and the order of
the stage markers.  Lemma~\ref{lem:completion-representation} extends
the dense representation to an ordinary countable representation of
$B\restr c$.  Full-super unembeddability of $I_\gamma$ into
$J_\gamma$ now allows Proposition~\ref{prop:exclusion} to be applied,
which excludes the asserted embedding.
\end{proof}

% End surgery.tex
% Begin construction.tex
\section{Construction of the counterexample}\label{sec:construction}

We now construct the Boolean retraction required by
\cref{prop:transfer}. Recall that $\lambda=(2^\cc)^+$.

\begin{theorem}\label{thm:boolean}
There are nonzero complete ccc Boolean algebras $D\subseteq B$, each of
cardinality $\lambda$, and a Boolean homomorphism
\[
 p:B\longrightarrow D,\qquad p|_D=\id_D,
\]
with the following properties:
\begin{enumerate}
\item The ideal $K=\ker p$ is order dense in $B$.
\item For every $0<b\in K$ and $0<d\in D$, there is no Boolean embedding
\[
 D\restr d\ \longrightarrow\ B\restr b.
\]
\item Every positive interval of $D$ has cardinality $\lambda$.
\end{enumerate}
\end{theorem}

The construction uses a fixed complete algebra $Q$ in which every positive
interval contains each $E_\gamma$. We build a complete ccc algebra $B$
and a homomorphism $q:B\to Q$, inserting $G_\gamma$ on the complement
of each positive element killed by $q$. The complementary-interval lemma
will then exclude $E_\gamma$ from that element's interval. The inserted
copies also provide lifts of the free factors of a dense subalgebra of
$Q$; extending these lifts gives a section of $q$. Its image is the
required complete subalgebra $D$.

\subsection{A fixed complete target}

We first recall the extension property of complete Boolean algebras.

\begin{lemma}[Boolean extension theorem]\label{lem:boolean-extension}
Let $L$ be a subalgebra of a Boolean algebra $R$, and let $Q$ be a
complete Boolean algebra. Every Boolean homomorphism $h:L\to Q$
extends to a Boolean homomorphism $R\to Q$.
\end{lemma}

\begin{proof}
For $b\in R$, choose $y\in Q$ satisfying
\[
 \bigvee\{h(a):a\in L,\ a\leq b\}
 \ \leq y\leq\
 \bigwedge\{h(a):a\in L,\ b\leq a\}.
\]
Such a choice exists because $a\leq b\leq a'$ implies
$h(a)\leq h(a')$. The formula
\[
 (a_1\wedge b)\vee(a_0\wedge\neg b)
 \longmapsto
 (h(a_1)\wedge y)\vee(h(a_0)\wedge\neg y)
\]
defines a homomorphism on the subalgebra $L[b]$. Indeed,
$a_1\wedge b=0$ implies $y\leq h(\neg a_1)$, and
$a_0\wedge\neg b=0$ implies $h(a_0)\leq y$.
Thus zero cells have zero image, which makes the formula independent of
the expression chosen; it then preserves the Boolean operations.
Zorn's lemma applied to these extensions, with codomain $Q$ fixed,
gives an extension to all of $R$.
\end{proof}

\begin{lemma}\label{lem:dense-fixing}
Let $H$ be a dense subalgebra of a Boolean algebra $Q$. A Boolean
endomorphism of $Q$ that fixes $H$ pointwise is the identity.
\end{lemma}

\begin{proof}
Let $h$ be such an endomorphism and $x\in Q$. If
$x\wedge\neg h(x)>0$, density supplies $0<a\in H$ below this
element. But $a=h(a)\leq h(x)$, a contradiction. Hence
$x\leq h(x)$. Applying the same argument to $\neg x$ gives the
opposite inequality.
\end{proof}

Recall the countable free products $G_\gamma$ from
\cref{lem:countable-free-product}. Form the ordinary Boolean free product
and its completion
\[
 H=*_{\gamma<\lambda}G_\gamma,
 \qquad Q=\widehat H.
\]
Write $\iota_\gamma:G_\gamma\hookrightarrow Q$ for the factor
inclusions, and write
$\iota_{\gamma,m}:E_\gamma\hookrightarrow Q$ for their restrictions
to the $m$th copies, using the chosen identifications
$E_\gamma^{(m)}\cong E_\gamma$.

\begin{lemma}\label{lem:universal-target-intervals}
For every $\gamma<\lambda$ and every $0<x\in Q$, there is a Boolean
embedding $E_\gamma\hookrightarrow Q\restr x$ preserving the
relative unit.
\end{lemma}

\begin{proof}
Choose $0<h\in H$ with $h\leq x$. A finite expression for $h$
uses only finitely many factors $E_\delta^{(m)}$. Choose $m<\omega$
such that the factor $E_\gamma^{(m)}$ is unused. Independence of the
free factors implies
\[
 h\wedge\iota_{\gamma,m}(e)>0
 \qquad(0<e\in E_\gamma).
\]
Choose an ultrafilter homomorphism $\varepsilon:E_\gamma\to2$.
The map
\[
 e\longmapsto
 \bigl(h\wedge\iota_{\gamma,m}(e)\bigr)
 \vee\varepsilon(e)(x\wedge\neg h)
\]
is a Boolean homomorphism with relative unit $x$. Its components have
disjoint units $h$ and $x\wedge\neg h$, and its first component is
injective by the displayed positivity. Here, as below,
$\epsilon a$ means $0$ or $a$ according as $\epsilon\in2$ is $0$ or $1$.
\end{proof}

\subsection{The source algebras and the map to \texorpdfstring{$Q$}{Q}}

We use the cardinal identity $\lambda^{\aleph_0}=\lambda$. Indeed,
$(2^\cc)^{\aleph_0}=2^\cc$, and every countable subset of the regular
cardinal $\lambda=(2^\cc)^+$ is bounded. There are $\lambda$ possible
bounds, each of cardinality at most $2^\cc$. Consequently a ccc Boolean
algebra of cardinality at most $\lambda$ has completion of cardinality
at most $\lambda$: each element of the completion is the supremum of a
countable antichain in the dense algebra.

For $\alpha\leq\lambda$, we construct complete ccc Boolean algebras
$C_\alpha$, dense subalgebras $F_\alpha$ with ordinary finitary
representations, and compatible homomorphisms
\[
 q_\alpha:C_\alpha\longrightarrow Q.
\]
The inclusions $C_\alpha\subseteq C_\beta$ will be complete, and
$q_\beta|_{C_\alpha}=q_\alpha$ whenever $\alpha<\beta$.
All the source algebras will have cardinality at most $\lambda$.
At successor stage $\gamma$ we also record a unital homomorphism
\begin{equation}\label{eq:recorded-lift}
 g_\gamma:G_\gamma\longrightarrow C_{\gamma+1},
 \qquad q_{\gamma+1}g_\gamma=\iota_\gamma.
\end{equation}

\paragraph{The initial algebra.}
Set $C_0=\mathcal P(\omega)$, and let $F_0$ be its finite--cofinite
subalgebra. Give every element of $F_0$ a distinct reserved nullary
code. This is an ordinary finitary representation, and $F_0$ is
countable and dense in $C_0$. Choose a free ultrafilter on $\omega$;
its homomorphism $C_0\to2$, followed by $2\hookrightarrow Q$, defines
$q_0$. For the singleton elements $a_n=\{n\}$ we have
\begin{equation}\label{eq:initial-kernel-density}
 q_0(a_n)=0\quad(n<\omega),
 \qquad \bigvee_{n<\omega}^{C_0}a_n=1.
\end{equation}
The complete source embeddings will preserve this join.

\paragraph{An assignment of stages.}
Fix an injection
\[
 \tau:\lambda\times\lambda\longrightarrow\lambda,
 \qquad \tau(\beta,\xi)\geq\beta.
\]
To construct it, enumerate $\lambda\times\lambda$ in order type
$\lambda$ and, for each pair $(\beta,\xi)$, choose an unused ordinal
at least $\beta$. Fewer than $\lambda$ ordinals have been used, while
the available tail has cardinality $\lambda$.

For each $\beta<\lambda$, once $C_\beta$ has been formed, enumerate
without repetitions the set
\[
 C_\beta\setminus\bigcup_{\eta<\beta}C_\eta
   =\{x_{\beta,\xi}:\xi<\theta_\beta\},
 \qquad\theta_\beta\leq\lambda.
\]
Assign $x_{\beta,\xi}$ to stage $\tau(\beta,\xi)$. At most one
element is assigned to each stage, and that element has already been
constructed when its assigned stage is reached.

\paragraph{Successor stages.}
Suppose $C_\alpha,F_\alpha,q_\alpha$ have been constructed. If the
element assigned to stage $\alpha$ is positive and has image zero
under $q_\alpha$, call it $b_\alpha$. Otherwise set $b_\alpha=0$.
Put $u_\alpha=\neg b_\alpha$. Since
$q_\alpha(u_\alpha)=1$, this unit is nonzero. Apply
\cref{lem:surgery} to insert $G_\alpha$ at $u_\alpha$, obtaining
\begin{equation}\label{eq:construction-surgery}
 C_{\alpha+1}
  =(C_\alpha\restr b_\alpha)\times
      \widehat{(C_\alpha\restr u_\alpha)*G_\alpha}.
\end{equation}
When $b_\alpha=0$, the zero factor is omitted. Thus there is an
insertion at every stage, including stages with no positive kernel
element assigned to them. The lemma supplies the dense algebra
$F_{\alpha+1}$ and its ordinary finitary representation.

To extend $q_\alpha$, consider the dense subalgebra
\[
 R_\alpha=(C_\alpha\restr b_\alpha)\times
             \bigl((C_\alpha\restr u_\alpha)*G_\alpha\bigr)
 \ \subseteq C_{\alpha+1}.
\]
The restriction $q_\alpha:C_\alpha\restr u_\alpha\to Q$ is
unital. Together with $\iota_\alpha:G_\alpha\to Q$, it induces a
homomorphism from the second free-product factor to $Q$. Composing
with the projection onto that factor gives
$\rho_\alpha:R_\alpha\to Q$. For $c\in C_\alpha$,
\[
 \rho_\alpha(c)=q_\alpha(c\wedge u_\alpha)=q_\alpha(c),
\]
because $q_\alpha(b_\alpha)=0$.
By \cref{lem:boolean-extension}, $\rho_\alpha$ extends to
$q_{\alpha+1}:C_{\alpha+1}\to Q$.

Let $\ell_\alpha:G_\alpha\to C_{\alpha+1}\restr u_\alpha$
be the fresh relative factor embedding. Choose an ultrafilter
homomorphism $\varepsilon_\alpha:G_\alpha\to2$, and define
\begin{equation}\label{eq:unital-splice}
 g_\alpha(e)=\ell_\alpha(e)\vee\varepsilon_\alpha(e)b_\alpha.
\end{equation}
The two components have complementary units, so $g_\alpha$ is a
unital Boolean homomorphism. Its image under $q_{\alpha+1}$ is
$\iota_\alpha(e)$, proving \eqref{eq:recorded-lift}.

The source codes use only the old source coefficients and the new
tree $I_\alpha$. Thus before stage $\alpha$ their tree indices are
strictly less than $\alpha$, as required by \cref{lem:surgery}.
The values of $q_\alpha$ impose no additional requirement on these
codes: the homomorphism out of the free product exists for any pair
of homomorphisms from its two factors. By
\cref{lem:countable-free-product}, $|G_\alpha|=\lambda$;
\cref{lem:surgery} and the completion bound above therefore give
$|C_{\alpha+1}|\leq\lambda$ and the required ccc and complete
embedding properties.

\paragraph{Limit stages.}
At a nonzero limit $\delta\leq\lambda$, set
\[
 F_\delta=\bigcup_{\alpha<\delta}F_\alpha,
 \qquad C_\delta=\widehat{F_\delta}.
\]
The continuous regular-chain assertion in \cref{lem:surgery} gives
the ccc, the retained ordinary representation, and complete embeddings
of all earlier $C_\alpha$ into $C_\delta$. In particular,
$\bigcup_{\alpha<\delta}C_\alpha$ is a subalgebra of $C_\delta$.
The compatible maps $q_\alpha$ define a homomorphism on this union;
\cref{lem:boolean-extension} extends it to $q_\delta:C_\delta\to Q$.
All previously recorded identities \eqref{eq:recorded-lift} persist.
The size of the dense union, followed by the completion bound, gives
$|C_\delta|\leq\lambda$. This completes the recursion.

\subsection{The section and the interval obstruction}\label{sec:retraction}

\begin{proof}[Proof of \cref{thm:boolean}]
Put $B=C_\lambda$ and $q=q_\lambda$. The recorded maps
$g_\gamma:G_\gamma\to B$ induce a homomorphism
\[
 j:H=*_{\gamma<\lambda}G_\gamma\longrightarrow B.
\]
By \eqref{eq:recorded-lift}, $qj:H\to Q$ is the canonical inclusion.
Since $B$ is complete, \cref{lem:boolean-extension} extends $j$ along
$H\subseteq Q$ to a homomorphism $s:Q\to B$. The endomorphism $qs$
of $Q$ fixes the dense subalgebra $H$, so \cref{lem:dense-fixing} gives
\begin{equation}\label{eq:constructed-section}
 qs=\id_Q.
\end{equation}
Consequently $s$ is injective. Define
\begin{equation}\label{eq:constructed-retraction}
 D=s(Q)\subseteq B,
 \qquad p=sq:B\longrightarrow D,
 \qquad K=\ker q=\ker p.
\end{equation}
Equation \eqref{eq:constructed-section} shows that $p$ fixes $D$.
The algebra $D$ is complete because it is isomorphic to $Q$, and it
is ccc as a subalgebra of $B$. Its joins are computed internally;
the embedding $s$ need not preserve arbitrary joins in $B$.

\paragraph{Density of the kernel.}
The original singleton elements $a_n$ still have supremum $1$ in
$B$ and are killed by $q$, by \eqref{eq:initial-kernel-density}.
For every $b>0$ in $B$, distributivity of a finite meet over this
join gives $b\wedge a_n>0$ for some $n$. This positive element lies
in $K$ below $b$. Thus $K$ is order dense.

\paragraph{Every positive kernel element was assigned a stage.}
Every $b\in B$ belongs to some $C_\beta$ with $\beta<\lambda$.
Indeed, choose a maximal antichain of the dense algebra $F_\lambda$
below $b$. It is countable, and its supremum in $B$ is $b$.
Its members lie in some $F_\beta$ because
$\operatorname{cf}(\lambda)>\aleph_0$. Their supremum in $C_\beta$
equals their supremum in $B$, since $C_\beta\subseteq B$ is complete.
This proves
\begin{equation}\label{eq:bounded-source}
 B=\bigcup_{\beta<\lambda}C_\beta.
\end{equation}
In particular, every element has a first stage of appearance and is
assigned a stage by the injection $\tau$.

Fix $0<b\in K$, and let $\gamma$ be its assigned stage. Compatibility
of the maps to the fixed algebra $Q$ gives $q_\gamma(b)=q(b)=0$.
Hence $b_\gamma=b$ and $G_\gamma$ was inserted at $\neg b$.
By \cref{lem:complementary-omission},
\begin{equation}\label{eq:kernel-tree-exclusion}
 E_\gamma\not\hookrightarrow B\restr b.
\end{equation}
On the other hand, \cref{lem:universal-target-intervals} and the
isomorphism $s:Q\xrightarrow{\sim}D$ give, for every $0<d\in D$,
\begin{equation}\label{eq:tree-in-target-interval}
 E_\gamma\hookrightarrow D\restr d.
\end{equation}
An embedding $D\restr d\hookrightarrow B\restr b$ would contradict
\eqref{eq:kernel-tree-exclusion} by composition. This proves the
required interval obstruction for every $b>0$ in $K$ and $d>0$ in $D$.

Finally, \cref{lem:universal-target-intervals} applies to any fixed
$\gamma<\lambda$, so $|E_\gamma|=\lambda$ gives
\[
 \lambda\leq |D\restr d|\leq |D|\leq |B|\leq\lambda
 \qquad(0<d\in D).
\]
Both algebras are nonzero since $Q$ and $B$ are nontrivial unital
Boolean algebras. These inequalities prove all the cardinality
assertions.
\end{proof}

\begin{proof}[Proof of \cref{thm:main}]
Apply \cref{prop:transfer} to the inclusion and retraction supplied by
\cref{thm:boolean}. Explicitly, set $S=\Stone(B)$ and $T=\Stone(D)$,
and let $r:S\to T$ and $i:T\to S$ be dual to the inclusion and the
retraction, respectively. Then $ri=\id_T$, and
\[
 P=\ker\bigl(\Z[r]:\Z[S]\to\Z[T]\bigr)
\]
is nonzero compact projective and is not isomorphic to $\Z[X]$ for any
compact Hausdorff extremally disconnected space $X$. The section
$\Z[i]$ gives the asserted splitting
$\Z[S]\cong\Z[T]\oplus P$.
\end{proof}

The same interval obstruction identifies the kernel from the inclusion
$D\subseteq B$ alone and determines its retraction uniquely.

\begin{corollary}\label{cor:intrinsic-kernel}
For the algebras in \cref{thm:boolean}, put
\[
 J_D(B)=\{b\in B:\text{no positive interval of }D
                       \text{ embeds into }B\restr b\}.
\]
Then $J_D(B)=\ker p$, and $p$ is the unique Boolean homomorphism
$B\to D$ whose restriction to $D$ is the identity.
\end{corollary}

\begin{proof}
The interval obstruction gives $\ker p\subseteq J_D(B)$, including
$b=0$. Conversely, suppose $p(b)=d_0>0$. Choose an ultrafilter
homomorphism $\varepsilon:D\restr d_0\to2$. The formula
\begin{equation}\label{eq:interval-section}
 D\restr d_0\longrightarrow B\restr b,
 \qquad x\longmapsto(b\wedge x)\vee
                  \varepsilon(x)(b\wedge\neg d_0)
\end{equation}
defines a Boolean homomorphism with relative unit $b$: its two
components have disjoint units $b\wedge d_0$ and $b\wedge\neg d_0$.
Applying $p$ to its value gives $x$. It is therefore injective, so
$b\notin J_D(B)$. Hence
\begin{equation}\label{eq:intrinsic-ideal}
 J_D(B)=\ker p.
\end{equation}

If $p':B\to D$ is another retraction, the same formula with $p'$
shows that $p'(b)>0$ implies $b\notin J_D(B)$. Thus $p'$ kills
$\ker p$. For every $b\in B$ we have
$b\mathbin{\triangle}p(b)\in\ker p$. Applying $p'$ and using
$p'|_D=\id_D$ gives $p'(b)=p(b)$, proving uniqueness.
\end{proof}

\begin{remark}[A smaller admissible cardinal]\label{rem:smaller-cardinal}
The entire Boolean construction, including its interval obstruction and
uniqueness conclusion, also works with
\[
 \lambda_*=(2^{\aleph_1})^+\leq\lambda.
\]
The inequality need not be strict.  The refinement uses a different coding
of term shapes in \cref{lem:three-successors}; the term algebra remains
$M_{\cc,\aleph_1}(J)$.

There are at most $\cc$ countable term shapes.  Assign to each shape
$\sigma$ a distinct binary sequence $b_\sigma\in2^\omega$.  For an
omitted-tree sum $J$ used in the construction, fix its root $o$ and a
nonroot index $t$. Such sums are nontrivial. For each term $f(v)$, enumerate its
index occurrences as $(a_{v,k})_{k<\omega}$, retaining repetitions and
padding unused positions by $o$; its shape records which positions are
used.  Replace the variable-length code in that lemma by
\[
 \begin{split}
 f_1(v)(2k)&=
 \begin{cases}
 o,&b_{\sigma_v}(k)=0,\\
 t,&b_{\sigma_v}(k)=1,
 \end{cases}\\
 f_1(v)(2k+1)&=a_{v,k},
 \end{split}
 \qquad k<\omega,
\]
where $\sigma_v$ is the shape of $f(v)$.  Include $t$, the coding choices,
and $f_1$ in the prescribed parameter for super-unembeddability.

Now $(\aleph_1,\aleph_1)$-super-unembeddability suffices: every code
sequence has length $\omega<\aleph_1$.  The three successor sequences
belong to $N_n$, and so do all their coordinates, since
$\omega\subseteq M_0$.  Equality of their atomic types recovers every
shape bit through the tests $z_{2k}=o$ and preserves the joint types of
the odd-coordinate index lists.  By \cref{lem:external-parameters}, these
types remain equal after adjoining the finite support of $f(\eta)$,
once $n$ is sufficiently large.  Thus the same three pairs are similar,
which proves the successor-similarity and completion-exclusion statements
for these target trees under the weaker super-unembeddability hypothesis.

The required full family exists by \cite[Lemma~1.11(1)]{Shelah2023}:
its hypotheses are that $\lambda_*$ is regular,
$\lambda_*>\mu\geq\kappa$, and $\lambda_*>\mu^\kappa$.
Take $\mu=\kappa=\aleph_1$ and use
$(\aleph_1)^{\aleph_1}=2^{\aleph_1}$.
Also $\lambda_*^{\aleph_0}=\lambda_*$: every countable sequence in
$\lambda_*$ is bounded, and
$(2^{\aleph_1})^{\aleph_0}=2^{\aleph_1}$.
Consequently the completion size bounds and bookkeeping remain valid.
With this family and the modified similarity argument, the represented
intervals, local insertions, fixed-target construction, uniqueness proof,
and passage to condensed abelian groups apply unchanged.
\end{remark}


% End construction.tex
% Begin references.tex
\begin{thebibliography}{9}
\raggedright

\bibitem{ClausenScholzeComplex}
D.~Clausen and P.~Scholze,
\emph{Condensed Mathematics and Complex Geometry},
arXiv:2605.11731v1, May 2026.
\url{https://arxiv.org/abs/2605.11731}.

\bibitem{ScholzeCondensed}
P.~Scholze,
\emph{Lectures on Condensed Mathematics},
lecture notes, May 2019; all results joint with D.~Clausen.
\url{https://www.math.uni-bonn.de/people/scholze/Condensed.pdf}.

\bibitem{Shelah2023}
S.~Shelah,
\emph{Building complicated index models and Boolean algebras},
Sh:511, version May 23, 2023.
\url{https://shelah.logic.at/files/230169/511.pdf}.

\bibitem{Shelah2026}
S.~Shelah,
\emph{Building complicated index models and Boolean algebras},
Sh:511, version June 9, 2026.
\url{https://shelah.logic.at/files/255135/511.pdf}.

\bibitem{ShelahE59}
S.~Shelah,
\emph{General non-structure theory},
Sh:E59, version October 11, 2024.
\url{https://shelah.logic.at/files/242052/E59.pdf}.

\end{thebibliography}

% End references.tex
\end{document}
